この度2025年度のアクチュアリー１次試験「数学」を受験したので、それの受験記を書きます。

<h2>受験に至る経緯</h2>
　以前からアクチュアリーという職業に興味はあったため、ふと試験の日程を確認したところ奇跡的に出願締切が明日10月8日だったので、なんとなく受けてみようと思い申し込みました。

<h2>10月にしたこと</h2>
　まず最初に使った参考書は『1冊でマスター 大学の統計学』です。この参考書は自分が以前から持ってたので使いました（内容はほぼ未読の状態）。最初の3日ほどで全体を軽く二周した程度でしか使ってないのと、後述の合ストにほぼ書いてあるので別にやらなくても良かったと思います。<br><br>

　次の参考書に取り組む前に、自分の今の現在地を確認するために2024年の過去問を解きました（結果は惨敗）。マジで知らない概念が多くて焦りつつも、案外数学系の問題は解けたので自信がついた反面、不安も残る結果となりました。<br><br>

　次に使った参考書は『アクチュアリー試験 合格へのストラテジー 数学』（以下、合スト）です。10月は公式などは特に覚えたりはせず、理解して使えるようになるというのを目標に一旦終わらせました。
アク試は「確率」「統計」「モデリング」の３分野に分かれているのですが、確率に関しては受験数学や競技数学で触れていた部分だったのであまり勉強はしませんでした。ただ、モデリング分野は自分が全く知らない概念だったのでかなり苦労しました。<br><br>
　10月の合計勉強時間は約50時間でした。

<h2>11月にしたこと</h2>
　11月頭に、自分で各分布の期待値・分散・確率密度関数や、統計量・信頼区間などを纏めた暗記プリントを作成し、数日に1回は解くようにしました。導出が簡単に出来るものもありますが、負の二項分布とかF分布などは導出が面倒くさかったり難しかったりするので、覚えておいて損は無いと思います。これのお陰で過去問演習や試験本番は困ることが無かったのでやって良かったと思っています。<br><br>
　11月は主に 過去問演習→解けなかった問題を合ストで確認する という作業をしました。過去問を遡っていくと、「これ合ストで見たやつだな～」とか「他の年度の問題にもあったな～」となることが、かなりあります。他のサイトでも散々書かれていることですが、それくらい過去問は大事です。過去問は2024年から平成23年まで解きました。<br><br>
　11月の勉強時間は約60時間でした。

<h2>試験直前！12月にしたこと</h2>
　これまではアプリの電卓を使っていたのですが、そろそろ慣れないといけないなと思い電卓を買いました（<a href="https://www.amazon.co.jp/dp/B0FTFH8Q99?tag=note0e2a-22&linkCode=ogi&th=1&psc=1">私が購入したものはこちら</a>）。
あくまで主観ですが、CMとRMが独立しているものを選ぶと楽な気がします。かなり使いやすかったのでオススメです。<br>
　12月は過去問を平成18年まで解きました。ただ、配点が平成23年から公開されていないのと、11月末時点で合格点＋αの点数を取れていたので、別にやらなくても良かったような気はします。<br>
　試験直前ということもあったので、基本的には過去問で間違えてた問題を詰め直したり合ストに戻って細かい知識の確認をしました。<br>
　12月の勉強時間は約30時間でした。


<h2>試験当日！（以降問題のネタバレを含みます）</h2>
　試験開始は当日の14:30からと、かなり時間に余裕があります。昼食を済ませ、コンビニで買ったカフェラテを飲んで待ってました。<br>
　初のCBT試験だったので少し早めに会場入りしてパラパラと合ストをめくり、他の受験者も会場入りして少し経ったとき、ふと周りを見ると全員が合ストを読んでいてびっくりしました。中には電卓３個持ちの人も居て猛者っぷりが凄かった。

<h3>試験開始！</h3>
問題は <a href="https://www.actuaries.jp/examin/info.html", target="_blank">こちら</a> から確認できます。<br><br>
　第１～３問は５点 ✕ 10問、計50点分の小問集合です。<br>
　私の構想では小問集合で45点、後半で30点稼ぐはずでしたが、第２問（２）の誤差の解法をド忘れするというハプニングが発生（過去問や合ストでも散々解いたのに…）。次ぐ（４）でも変動係数とかいうもはや誰も覚えていないやつが出題され敗北…。既に滅茶苦茶テンションが下がっているが、第３問ではボーナスとも言える回帰式とマルコフ連鎖が出題。これのお陰で気を取り戻し、リフレッシュのために一旦トイレへ。<br><br>

　第４問はシミュレーション。シミュレーションは解法さえ知っていれば誘導に乗るだけのパターンが多いので着々と手を進める。順調に（１）を解き終え（２）の積分に取り掛かる。<br><br><br>

……合わない。何度計算しても⑧の計算結果が解答欄に無い。汗が滝のように溢れ出る。シミュレーションの解法を間違って覚えていた…？そうなると（１）も間違っている可能性が高い。３回ほど計算しても合わず、しかも積分がまじで面倒くさい。ここで不合格の３文字が脳裏によぎる。
時間を無駄にしないために第４問を捨てた。もしかしたら（３）は解けたかもしれないが、そんなことを考えている余裕など微塵もなかった。<br><br>

　続いて第５問を見る。……条件付き期待値？？？？？<br>
過去問の数学問題はほぼ解けていたが、唯一苦手だった条件付き期待値。
試験数日前に条件付き確率が出題されなければ勝てると豪語していたものの、出題されてしまう。ここで不合格を確信する。
ただ全く手をつけないというわけもいかず、一競技数学徒としてなんとか（２）まで解いたものの（３）で新たな添字が登場し、脳がバグる前に諦めることにした。<br><br>

　最後に第６問。テーマはポアソン分布の精密法の導出。
ここにきて救われる。平成22年の過去問をやっていてよかった。ありがとう12月の俺。難なく突破する。<br><br>

　残り時間、第４問と第５問に立てていた大量のフラグを勘で埋めて終了。手応えは全く無く、ただ生きた心地がしなかった。<br><br>

<h2>結果＆感想</h2>
　結果発表は２ヶ月後の２月中旬。受験後に何度か自己採点をしてみたら意外と点数は悪すぎるということはなく、受かってるか落ちているか五分五分の状態だったので、最後に勘で塗ったヤツでなんとか受かってくれ～と思ってました。<br><br>

結果は...<br><br><br>

<img src="/articles/pages/001_gokaku.jpg"><br>
まさかの合格！！流石に脳汁出まくって感謝の正拳突きしました。<br><br>
　10月になんとなくで申し込んでから約２ヶ月、計140時間の勉強で無事合格することができました。何より、『知らないことを学ぶ楽しさ』を再確認できた非常に濃密な時間でした。<br><br>



　最後に、ここまで読んでいただきありがとうございます。この記事を読んでアクチュアリー試験に興味を持っていただけたら嬉しいです。折角なので、私の足跡として過去問演習の点数推移を公開しておきます。これから挑む方の参考になれば嬉しいです（→<a href="https://yura1685.github.io/articles/pages/001_actuary_math_score.pdf" target="_blank">過去問の点数</a>）。<br><br>

それでは、次の科目でお会いしましょう！